\chapter{Stand der Forschung}
\label{cha:Stand der Forschung}

\section{Autonomes Landen von Drohnen}
\label{sec:landen}

Bisherige Arbeiten zum autonomen Landen von Drohnen lassen sich in zwei grundlegende Kategorien einteilen: traditionelle und bildbasierte Methoden.
Traditionelle Methoden stützen sich auf Signale externer Infrastruktur wie GNSS-Satelliten oder Funkbaken, um Position und Lage der Drohne zu bestimmen.
Bildbasierte Methoden hingegen nutzen Kameradaten zur Wahrnehmung der Umgebung und lassen sich weiter unterteilen in markerbasierte Systeme, die definierte
visuelle Referenzpunkte zur Lokalisierung verwenden, und markerfreie Systeme, die adäquate Landezonen eigenständig aus Bilddaten ableiten und somit keine
zusätzliche Bodeninfrastruktur erfordern. Diese Einteilung orientiert sich an der Taxonomie von \textcite{arafatVisionBasedNavigationTechniques2023}.

Orthogonal zur Sensorik lassen sich Ansätze danach unterscheiden, welche Teilaufgaben des Guidance, Navigation and Control (GNC) \autocite{kanellakisSurveyComputerVision2017} sie adressieren und wie diese zusammenwirken. Modulare Methoden zerlegen das Problem in getrennte Komponenten wie Obstacle Detection, Pose-Estimation oder Path-Planning, die als aufeinanderfolgende Stufen einer Pipeline gelöst werden \autocite{kanellakisSurveyComputerVision2017, wangVisionBasedDeepReinforcement2024}. End-to-End-Methoden hingegen bilden Sensorinput direkt auf Steuerkommandos ab, ohne solche expliziten Zwischenstufen \autocite{wangVisionBasedDeepReinforcement2024}.

GNSS ist die verbreitetste Methode zur Positionsbestimmung von Drohnen: Ein einzelnes GNSS-Modul liefert eine Genauigkeit von 5–10\,m und bildet die Grundlage für Missionsplanung und -ausführung \autocite{jarrayaGnssdeniedUnmannedAerial2025}. Dass autonomes Landen bei bekannter Relativposition grundsätzlich beherrschbar ist, zeigen \textcite{voosNonlinearTrackingLanding2010} mittels nichtlinearer Regelung auf stationären und beweglichen Zielen. Die zentrale Herausforderung liegt somit nicht im Regelungsentwurf, sondern in der Genauigkeit der zugrunde liegenden Positionsdaten.

In der Praxis erreicht GNSS-basierte Landung jedoch nur Meter-Genauigkeit. \textcite{patrikGNSSbasedNavigationSystems2019} berichten von einer durchschnittlichen Landeabweichung von 1,11\,m; auch beim Open-Source-Autopiloten PX4 beträgt die GNSS-basierte Landegenauigkeit mehrere Meter \autocite{meierPX4NodebasedMultithreaded2015, wubbenAccurateLandingUnmanned2019}, sodass für höhere Präzision zusätzliche Sensoren erforderlich sind \autocite{PrecisionLanding}. Die Positionsgenauigkeit lässt sich zwar durch RTK-Korrekturen auf Zentimeter-Niveau steigern — \textcite{tavasciReliabilityRealTimeKinematic2024} berichten zentimetergenaue Echtzeit-Navigation unter idealen Bedingungen —, allerdings fällt die Genauigkeit in schwierigen GNSS-Umgebungen, etwa bei Signalabschattung durch Vegetation oder Gebäude, auf Meter-Niveau zurück. Gerade in Agrarumgebungen mit Baumkronen und Laubdächern sind solche Bedingungen typisch.

GNSS bleibt damit als robuste Grobpositionierung unverzichtbar, weist jedoch für Präzisionslandung in strukturierten Umgebungen fundamentale Limitationen auf: Die Genauigkeit reicht nicht für das Landen zwischen Baumreihen oder Weinreben mit einem typischen Reihenabstand von 1,5–2\,m. Zudem liefert GNSS keinerlei Umgebungswahrnehmung; Veränderungen der Umgebung, Hindernisse oder beengte Platzverhältnisse bleiben unerkannt.

Bildbasierte Landesysteme mit visuellen Markern bieten eine deutlich höhere Präzision als rein GNSS-gestützte Methoden. So erreicht das ArUco-basierte System von \textcite{wubbenAccurateLandingUnmanned2019} durch deterministische Bildverarbeitung eine durchschnittliche Landeabweichung von 11\,cm bei zuverlässiger Markerdetektion aus bis zu 30\,m Höhe. Solche klassischen Ansätze setzen jedoch voraus, dass die Markererkennung unter allen Betriebsbedingungen robust funktioniert — eine Annahme, die bei Verdeckung, wechselnden Lichtverhältnissen oder beschädigten Markern nicht gegeben ist.

Reinforcement Learning bietet hier einen konzeptionellen Vorteil: Statt auf handspezifizierte Detektionsalgorithmen angewiesen zu sein, lernt ein RL-Agent eine End-to-End-Steuerung direkt aus Bilddaten und kann dabei implizit Robustheit gegenüber visuellen Störungen erwerben. \textcite{polvaraSimtoRealQuadrotorLanding2020} demonstrieren diesen Ansatz erstmals für die autonome Landung: Ihr System nutzt sequentielle Deep-Q-Networks mit diskretem Aktionsraum und einer Divide-and-Conquer-Strategie, die den Landeprozess in Marker-Ausrichtung und vertikalen Abstieg unterteilt — ein Ansatz, der das Problem dünn besetzter Belohnungssignale adressiert, indem die komplexe Aufgabe in einfachere Teilziele zerlegt wird. Als einziger Sensor dient eine niedrig aufgelöste monokulare Graustufenkamera. Ausschließlich in Simulation mit Domain Randomization trainiert, erzielt das System Erfolgsraten von 89\,\% für den vertikalen Abstieg, die im realen Flug auf 61\,\% sinken. Der Robustheitsvorteil gegenüber klassischen Methoden zeigt sich unter erschwerten Bedingungen: Bei verrauschten Markertexturen erreicht das RL-System noch 51\,\%, während ein konventioneller AR-Tracker vollständig versagt — allerdings ist auch dieser Wert für sicherheitskritische Anwendungen unzureichend. Eine zentrale methodische Einschränkung bleibt der diskrete Aktionsraum, der die Steuerungspräzision limitiert.

\textcite{ladoszAutonomousLandingMoving2024} erweitern den Ansatz auf bewegliche Landeplattformen und adressieren diese Einschränkung durch einen kontinuierlichen Aktionsraum, der erstmals auch die Vertikalgeschwindigkeit einschließt — bei früheren Arbeiten wird der vertikale Abstieg typischerweise regelbasiert gesteuert. Der Agent erhält neben Graustufenbildern auch IMU-Daten (Geschwindigkeit und Lage) als Zustandsinformation. In der Simulation erreicht das System Erfolgsraten von bis zu 97,4\,\% bei einer Landeabweichung von 0,52\,m; im realen Transfer sinkt die Rate auf rund 70\,\%. Die Autoren zeigen zudem, dass Domain Randomization zwar notwendig für den Sim-to-Real-Transfer ist, eine zu starke Randomisierung die Performance jedoch signifikant verschlechtert. Die Experimente beschränken sich auf einfache Szenarien ohne Hindernisse und gleichmäßig beleuchtete Innenräume.

Trotz dieser Fortschritte weisen markerbasierte Systeme grundlegende Einschränkungen auf, die unabhängig von der gewählten Steuerungsmethode gelten: Sie setzen vorplatzierte visuelle Infrastruktur am Landeziel voraus, was ihren Einsatz auf vordefinierte Orte begrenzt und in unbekannten Umgebungen unpraktikabel macht. Hinzu kommen methodische Limitationen der RL-basierten Ansätze: Beide Systeme wurden in vereinfachten Simulationsumgebungen ohne Hindernisse trainiert, und der Transfer in die reale Welt geht mit deutlichen Leistungseinbußen einher — \textcite{ladoszAutonomousLandingMoving2024} berichten einen Rückgang der Erfolgsrate von 97,4\,\% auf rund 70\,\% \autocite[vgl. auch][]{polvaraSimtoRealQuadrotorLanding2020}. Diese Limitationen motivieren die Entwicklung markerfreier Systeme, die Landezonen eigenständig aus Bilddaten ableiten und keine zusätzliche Bodeninfrastruktur erfordern.

Markerfreie Verfahren ermöglichen den Einsatz in unbekannten, unstrukturierten und dynamischen Umgebungen, da sie keine manuelle Vorbereitung der Landezone erfordern — eine Eigenschaft, die insbesondere für Notlandungen oder den Betrieb in unzugänglichen Gebieten entscheidend ist \autocite{yangMonocularVisionSLAMBased2018, bartolomeiAutonomousEmergencyLanding2022}. Die Bandbreite der Ansätze reicht dabei von klassischen modularen Pipelines bis hin zu End-to-End-Steuerungen mittels Deep Reinforcement Learning.

\textcite{yangMonocularVisionSLAMBased2018} verfolgen einen modularen Ansatz, bei dem aus monokularer SLAM-Rekonstruktion eine 3D-Punktwolke erzeugt wird, aus der anhand von Oberflächensemantik und Geländehöhe sichere Landezonen in einer Grid-Map identifiziert werden. Das System nutzt ausschließlich Visual-Inertial-Odometrie zur Zustandsschätzung und benötigt weder GNSS noch externe Lokalisierungsinfrastruktur. Es demonstriert Landungen in verschiedenen Umgebungen, weist jedoch aufgrund der hohen Pipeline-Komplexität eine Landezeit von rund zwei Minuten auf.

Einen End-to-End-Ansatz wählen \textcite{bartolomeiAutonomousEmergencyLanding2022}: Ihr RL-Agent bildet Tiefenbilder und semantische Segmentierungen — sogenannte Mid-Level-Repräsentationen, die generischer als Rohbilder sind und schnelleres Training ermöglichen — direkt auf diskrete Steuerkommandos ab. Das System erreicht in der Simulation Erfolgsraten von über 70\,\% bei minimaler Sensorausstattung mit lediglich einer monokularen RGB-Kamera. Ein methodisch relevanter Aspekt ist die Sim-to-Real-Strategie: Die Policy wird zunächst mit Ground-Truth-Tiefen- und Semantikdaten trainiert und anschließend mit verrauschten Prädiktionen neuronaler Netze feinabgestimmt, um den Simulationstransfer zu ermöglichen. Im Realflug landete die Policy erfolgreich, während sowohl eine kommerzielle als auch eine State-of-the-Art-Lösung am verrauschten Sensorinput scheiterten. Einschränkend setzt das System voraus, dass semantische Segmentierungen und Tiefenbilder stets verfügbar sind, und der diskrete Aktionsraum (zehn mögliche Aktionen) limitiert die Steuerungsflexibilität.

Dass Deep Reinforcement Learning auch unter dynamischen Umgebungsbedingungen effektive Landesteuerungen ermöglicht, zeigen \textcite{peterTornadoDroneBioinspiredDRLbased2024} mit TornadoDrone: Das bio-inspirierte DRL-Framework übertrifft eine konventionelle PID+EKF-Baseline bei Landungen unter starken Windturbulenzen deutlich, stützt sich in den realen Experimenten jedoch auf ein Vicon-Lokalisierungssystem, das unter Einsatzbedingungen nicht verfügbar ist.

\section{Tiefenbildbasierte Hindernisvermeidung}
\label{sec:tiefenbild-oa}

\textcite{loquercioLearningHighspeedFlight2021} demonstrieren ein End-to-End-Navigationssystem für Hochgeschwindigkeitsflüge durch komplexe, unbekannte Umgebungen mit ausschließlich bordeigener Sensorik und Rechenleistung. Das System nutzt Tiefenbilder einer Stereokamera und IMU-Daten, um kollisionsfreie Kurzzeittrajektorien zu erzeugen, die von einem nachgeschalteten PID-Regler ausgeführt werden. Trainiert wird ausschließlich in Simulation mittels Privileged Learning, einer Imitation eines Experten mit Zugang zu perfekter Zustandsinformation, wobei lediglich einfache Hindernisgeometrien (Zylinder, Quader, Baumstämme) verwendet werden. Durch die Simulation realistischen Sensorrauschens gelingt ein Zero-Shot-Transfer in reale Umgebungen, die während des Trainings nie erfahren wurden: dichte Wälder, verschneites Gelände und eingestürzte Gebäude. Bei niedrigen Geschwindigkeiten erreicht das System eine Erfolgsrate von 100\,\% in allen getesteten Umgebungen; im Vergleich zu modularen State-of-the-Art-Planern reduziert der Ansatz die Ausfallrate um den Faktor~10. Eine Einschränkung ist die geringe Positioniergenauigkeit, da Erfolg bereits bei einem Abstand von 5\,m zum Ziel definiert wird und der Fokus auf Geschwindigkeit statt Präzision liegt. Methodisch relevant für die vorliegende Arbeit ist der Nachweis, dass ein lernbasierter Ansatz mit Tiefenbildern und hierarchischer Architektur (Policy $\rightarrow$ PID) in dicht bewachsenen natürlichen Umgebungen robust navigieren kann. Die vorliegende Arbeit verfolgt eine vergleichbare Architektur, ersetzt jedoch Privileged Learning durch Reinforcement Learning und verschiebt den Fokus von Hochgeschwindigkeitsnavigation auf Präzisionslandung in beengten Umgebungen.

\textcite{kimMonocularVisionbasedAutonomous2022} verfolgen einen zweistufigen Ansatz, bei dem zunächst ein vortrainiertes CNN Tiefenbilder aus monokularen RGB-Aufnahmen schätzt und anschließend ein Dueling-Double-DQN auf Basis dieser geschätzten Tiefeninformation diskrete Ausweichentscheidungen trifft. Als einziger externer Sensor dient die bordeigene Monokularkamera der Drohne; Zustandsinformation wie Position oder Geschwindigkeit wird nicht benötigt. Das System wird in Simulation trainiert und ohne weitere Anpassung auf eine reale Drohne übertragen, wo es enge Korridore mit texturlosen Wänden, Personen und Kisten kollisionsfrei durchfliegt. Evaluiert wird anhand der zurückgelegten Strecke ohne Kollision; eine gezielte Navigation zu einem definierten Ziel ist nicht Teil des Systems. Die Testumgebungen beschränken sich auf strukturierte Innenräume; Experimente in natürlichen oder vegetationsreichen Umgebungen werden nicht durchgeführt. Der Ansatz demonstriert, dass selbst aus monokularer RGB-Eingabe geschätzte Tiefeninformation für lernbasierte Hindernisvermeidung ausreicht, allerdings mit den Einschränkungen eines diskreten Aktionsraums und der Abhängigkeit von der Qualität der Tiefenschätzung.

\textcite{xuNavRLLearningSafe2024} nutzen direkte Tiefenbilder, um mittels PPO eine kontinuierliche Policy für sichere Navigation in Umgebungen mit statischen und dynamischen Hindernissen zu trainieren. Die Architektur verarbeitet die Tiefenbilder zu einer 3D-Belegungskarte, die zusammen mit Drohnenzustandsdaten (Position, Geschwindigkeit, Orientierung) als Eingabe für das Actor-Critic-Netzwerk dient. Das Training erfolgt parallelisiert in Simulation mit Curriculum Learning, bei dem die Anzahl dynamischer Hindernisse schrittweise erhöht wird. Ein nachgeschalteter Safety Shield auf Basis von Velocity Obstacles ergänzt die gelernte Policy zur Laufzeit. Das System erreicht Zero-Shot-Transfer auf eine reale Drohne und erzielt in Szenarien mit dynamischen Hindernissen die geringste Kollisionsrate im Vergleich zu bestehenden Methoden. Auch hier beschränken sich die Testumgebungen auf strukturierte Indoor-Szenarien mit geometrisch einfachen Hindernissen.

Gemeinsam zeigen diese Arbeiten, dass tiefenbildbasierte lernende Systeme Hindernisse zuverlässig vermeiden können, unabhängig davon ob die Tiefeninformation direkt gemessen \autocite{loquercioLearningHighspeedFlight2021, xuNavRLLearningSafe2024} oder aus RGB-Bildern geschätzt wird \autocite{kimMonocularVisionbasedAutonomous2022}. Allerdings testen alle genannten Systeme in strukturierten Innenräumen oder Waldumgebungen; Evaluationen in landwirtschaftlichen Umgebungen mit Rebzeilen, niedrigen Baumkronen oder dichter Bodenvegetation fehlen.

Vereinzelte Arbeiten adressieren die Navigation in landwirtschaftlichen Umgebungen gezielt. \textcite{weiVisionBasedNavigation2025} demonstrieren ein lernbasiertes UAV-System für die autonome Navigation durch Obstplantagen-Reihen: Ein VAE-Controller, trainiert mittels Imitation Learning aus Expertendemonstrationen, steuert einen Quadrotor auf Basis einer front-montierten RGB-Kamera kollisionsfrei zwischen Baumreihen. Das System generalisiert auf ungesehene Umgebungen und kommt ohne GPS aus, beschränkt sich jedoch auf horizontale Reihennavigation. Für Weinberge zeigen \textcite{martiniPositionAgnosticAutonomous2022}, dass ein DRL-Agent auf Basis von Tiefenbildern ohne Positionsinformation durch simulierte Rebzeilen navigieren kann, allerdings mit einem Bodenroboter statt einer Drohne. Beide Arbeiten bestätigen, dass Vegetation in landwirtschaftlichen Umgebungen eine eigenständige Navigationsherausforderung darstellt, die durch lernbasierte Ansätze mit minimaler Sensorik adressiert werden kann. Keine der Arbeiten behandelt jedoch die vertikale Landung zwischen Vegetationsstrukturen, die eine kombinierte Hindernisvermeidung und Präzisionspositionierung erfordert.

\section{Reinforcement Learning für Drohnensteuerung}
\label{sec:architektur}

Über die Wahl der Sensorik hinaus beeinflusst die Steuerungsarchitektur die Leistungsfähigkeit lernbasierter Drohnensysteme maßgeblich, insbesondere der Aktionsraum und die Abstraktionsebene der Regelung. Die in den Abschnitten~\ref{sec:landen} und~\ref{sec:tiefenbild-oa} betrachteten Arbeiten lassen sich auch entlang dieser Dimension einordnen: End-to-End-Systeme wie \autocite{polvaraSimtoRealQuadrotorLanding2020, kimMonocularVisionbasedAutonomous2022, ladoszAutonomousLandingMoving2024} erzeugen Bewegungskommandos direkt aus Sensorinput, während hierarchische Ansätze wie \autocite{loquercioLearningHighspeedFlight2021, bartolomeiAutonomousEmergencyLanding2022} die Policy-Ausgabe an einen nachgeschalteten Regler delegieren. \textcite{xuNavRLLearningSafe2024} kombinieren beide Paradigmen, indem sie eine PPO-Policy mit einem regelbasierten Safety Shield ergänzen.

Die Wahl der Abstraktionsebene beeinflusst die Lernperformance substanziell: \textcite{esserActionSpaceDesign2024} zeigen in einer systematischen Studie über verschiedene Roboterplattformen und Aufgaben, dass der optimale Aktionsraum stark aufgabenabhängig ist und unterschiedliche Abstraktionsebenen zu deutlich verschiedenem Explorationsverhalten und finaler Performance führen. Speziell für Quadrotoren vergleichen \textcite{kaufmannBenchmarkComparisonLearned2022} drei Abstraktionsebenen: Geschwindigkeitskommandos, kollektiven Schub mit Drehraten sowie direkte Einzelmotorschübe. Die mittlere Ebene (Schub und Drehraten) erweist sich als beste Balance zwischen Agilität und Transferrobustheit; die höchste Abstraktion (Geschwindigkeit) schränkt die Manövrierfähigkeit ein, während die niedrigste (Einzelmotorschub) zu latenzempfindlich für den Realtransfer ist.

Neben der Steuerungsarchitektur hat sich PPO als verbreiteter Trainingsalgorithmus für RL-basierte Drohnensteuerung etabliert. Als Policy-Gradient-Verfahren eignet sich PPO für die kontinuierlichen Aktionsräume der Drohnensteuerung (vgl.\ Abschnitt~\ref{subsec:ppo_clip}). Da Reinforcement Learning grundsätzlich große Datenmengen benötigt, hat sich GPU-parallele Simulation als effektiver Ansatz etabliert: \textcite{rudinLearningWalkMinutes2022} zeigen, dass massiv-parallele Simulation mit hunderten bis tausenden Umgebungsinstanzen die Trainingszeit drastisch reduziert und schnellere Konvergenz ermöglicht. PPO eignet sich als On-Policy-Verfahren für dieses Paradigma, da es die parallel erzeugten Daten direkt verarbeitet und keinen Replay Buffer benötigt. Obwohl Off-Policy-Verfahren wie SAC eine höhere Sample-Effizienz bieten, zeigen \textcite{tayarAutonomousUAVFlight2025}, dass PPO bei sicherheitskritischen Präzisionsaufgaben in beengten Räumen zuverlässiger konvergiert als SAC. In den in diesem Kapitel betrachteten Arbeiten setzen \textcite{bartolomeiAutonomousEmergencyLanding2022}, \textcite{xuNavRLLearningSafe2024} und \textcite{kaufmannBenchmarkComparisonLearned2022} PPO ein; \textcite{kaufmannChampionlevelDroneRacing2023} demonstrieren den bislang eindrucksvollsten Einsatz, bei dem ein ausschließlich in Simulation trainiertes System drei menschliche Weltmeister im physischen Drohnenrennen schlägt.

Die vorliegende Arbeit folgt dem hierarchischen Paradigma und wählt eine noch höhere Abstraktionsebene als die von \textcite{kaufmannBenchmarkComparisonLearned2022} untersuchten: Der RL-Agent gibt eine Zielposition vor, die ein kaskadierender PID-Regler (Position $\rightarrow$ Geschwindigkeit $\rightarrow$ Lage $\rightarrow$ Drehzahl) in Motorkommandos umsetzt. Diese Wahl opfert bewusst Agilität zugunsten reduzierter Lernkomplexität, da die Aufgabe Präzisionslandung statt agilen Flugs erfordert. Die Ausgabe des Agenten bleibt interpretierbar, der PID-Regler kann unabhängig vom RL-Training auf die Zielplattform abgestimmt werden, und die Lernkomplexität reduziert sich auf die Frage, \emph{wohin} die Drohne fliegen soll.

Die Wahl der Sensorausstattung ist ein weiterer zentraler Designparameter: Drohnen sind aufgrund ihrer Energie- und Gewichtsbeschränkungen typischerweise nur mit minimaler Sensorik ausgestattet, bestehend aus Kameras und Distanzsensoren \autocite{bartolomeiAutonomousEmergencyLanding2022}. \textcite{semerikovVisionBasedAutonomousUAV2025} zeigen in ihrer Übersichtsarbeit, dass Vision-Inertial-Kombinationen eine gute Balance zwischen Wahrnehmungsleistung und Systemkomplexität bieten, während leistungsfähigere Sensorsetups entsprechend größere Drohnen mit höherer Lade- und Batteriekapazität erfordern. Gleichzeitig ermöglichen moderne Algorithmen, insbesondere lernbasierte Ansätze, dass auch mit leichtgewichtiger Sensorik anspruchsvolle Aufgaben gelöst werden können, die bisher komplexere Hardware voraussetzten. Die in diesem Kapitel betrachteten Arbeiten bestätigen diesen Trend: Erfolgreiche Hindernisvermeidung und Navigation gelingen mit einer einzelnen Kamera \autocite{kimMonocularVisionbasedAutonomous2022, polvaraSimtoRealQuadrotorLanding2020, bartolomeiAutonomousEmergencyLanding2022} oder einer Tiefenkamera mit IMU-Zustandsdaten \autocite{loquercioLearningHighspeedFlight2021, xuNavRLLearningSafe2024}. Die vorliegende Arbeit folgt diesem Prinzip minimaler Sensorik: Als Eingabe dienen eine nach unten gerichtete Tiefenkamera und ein Zustandsvektor (Position, Orientierung, Geschwindigkeit), wobei die algorithmische Komplexität des RL-Agenten die begrenzte Sensorausstattung kompensiert.

Tabelle~\ref{tab:vergleich-sdf} stellt die in diesem Kapitel betrachteten Arbeiten anhand ihrer Kernmerkmale gegenüber.

\begin{table}[htbp]
\centering
\caption{Vergleich der betrachteten Arbeiten zur autonomen Drohnennavigation und -landung. \emph{Hind.}: aktive Hindernisvermeidung; \emph{Landung}: Präzisionslandung als Teil der Aufgabe.}
\label{tab:vergleich-sdf}
\footnotesize
\setlength{\tabcolsep}{4pt}
\begin{tabular}{@{}lllcccc@{}}
\toprule
Arbeit & Sensorik & Methode & Akt.-raum & Hind. & Umgebung & Landung \\
\midrule
\textcite{polvaraSimtoRealQuadrotorLanding2020} & Mono RGB & DQN & diskret & -- & Indoor & $\checkmark$ \\
\textcite{ladoszAutonomousLandingMoving2024} & Mono RGB, IMU & DrQv2 & kont. & -- & Indoor & $\checkmark$ \\
\textcite{yangMonocularVisionSLAMBased2018} & Mono RGB & SLAM & -- & $\checkmark$ & Outdoor & $\checkmark$ \\
\textcite{bartolomeiAutonomousEmergencyLanding2022} & Mono RGB\textsuperscript{a} & PPO & diskret & $\checkmark$ & Outdoor & $\checkmark$ \\
\textcite{peterTornadoDroneBioinspiredDRLbased2024} & Vicon, IMU & TD3 & kont. & -- & Indoor & $\checkmark$ \\
\addlinespace
\textcite{loquercioLearningHighspeedFlight2021} & Stereo Tiefe, IMU & IL & kont. & $\checkmark$ & Outdoor & -- \\
\textcite{kimMonocularVisionbasedAutonomous2022} & Mono RGB\textsuperscript{b} & DQN & diskret & $\checkmark$ & Indoor & -- \\
\textcite{xuNavRLLearningSafe2024} & Tiefe, Zustand & PPO & kont. & $\checkmark$ & Indoor & -- \\
\textcite{weiVisionBasedNavigation2025} & Mono RGB & IL & kont. & $\checkmark$ & Agrar & -- \\
\textcite{martiniPositionAgnosticAutonomous2022}\textsuperscript{c} & Tiefe & SAC & kont. & $\checkmark$ & Agrar & -- \\
\addlinespace
Eigene Arbeit & Tiefe, Zustand & PPO & kont. & $\checkmark$ & Agrar & $\checkmark$ \\
\bottomrule
\end{tabular}

\smallskip
{\scriptsize
\textsuperscript{a} Eingabe: abgeleitete Tiefen- und Segmentierungsbilder. \quad
\textsuperscript{b} Eingabe: aus RGB geschätzte Tiefenbilder. \quad
\textsuperscript{c} Bodenroboter (UGV).
}
\end{table}

Zusammenfassend zeigt der Stand der Forschung drei wiederkehrende Limitationen RL-basierter Landesysteme: Erstens operieren die meisten Ansätze mit diskreten Aktionsräumen \autocite{polvaraSimtoRealQuadrotorLanding2020, bartolomeiAutonomousEmergencyLanding2022} oder vereinfachen den vertikalen Abstieg auf regelbasierte Steuerung — lediglich \textcite{ladoszAutonomousLandingMoving2024} setzen einen kontinuierlichen Aktionsraum mit Vertikalkontrolle ein, beschränken sich jedoch auf hindernisfreie Innenräume. Zweitens werden Hindernisvermeidung und Präzisionslandung überwiegend getrennt adressiert (vgl.\ Tabelle~\ref{tab:vergleich-sdf}): Systeme zur Hindernisvermeidung navigieren durch komplexe Umgebungen, ohne gezielt zu landen \autocite{loquercioLearningHighspeedFlight2021, kimMonocularVisionbasedAutonomous2022, xuNavRLLearningSafe2024}, während Landesysteme in der Regel in hindernisfreien Szenarien operieren \autocite{polvaraSimtoRealQuadrotorLanding2020, ladoszAutonomousLandingMoving2024}. Drittens bleibt der Sim-to-Real-Transfer eine offene Herausforderung: Die Erfolgsraten sinken systematisch beim Übergang in die reale Welt — von 89\,\% auf 61\,\% \autocite{polvaraSimtoRealQuadrotorLanding2020}, von 97,4\,\% auf rund 70\,\% \autocite{ladoszAutonomousLandingMoving2024} —, und reale Experimente stützen sich auf Hilfsmittel wie Vicon-Lokalisierung \autocite{peterTornadoDroneBioinspiredDRLbased2024} oder gezieltes Fine-Tuning \autocite{bartolomeiAutonomousEmergencyLanding2022}, die unter Einsatzbedingungen nicht verfügbar sind.

Die vorliegende Arbeit adressiert die ersten beiden Limitationen: Sie untersucht, ob ein RL-Agent mit kontinuierlichem Aktionsraum und Tiefenbildern in Simulation lernen kann, in beengten Umgebungen mit Hindernissen autonom zu landen. Ein vollständiger Sim-to-Real-Transfer liegt außerhalb des Rahmens dieser Arbeit. Um die Transferfähigkeit der gelernten Strategie dennoch zu untersuchen, wird ein Zero-Shot-Transfer in eine geometrisch verschiedenartige Simulationsumgebung (Sim-to-Sim) als vereinfachte Annäherung an die dritte Limitation eingesetzt.
